\section*{Chapitre \ref{ch:g7:how-animals-breathe} --- Comment les animaux respirent}

\begin{solution}{exo:g7:how-animals-breathe:1}
Une surface d'échange grande, mince et humide où le \omterm{prop:g4:journey-of-food:absorb}{sang} (ou le
liquide interne) rencontre le milieu. Grande~: assez de passage au
total~; mince~: un passage facile~; humide~: les gaz traversent
dissous.
\end{solution}

\begin{solution}{exo:g7:how-animals-breathe:2}
\omterm{def:g7:how-animals-breathe:four}{Poumons}~: \omterm{prop:g3:sorting-living-things:groups}{mammifères}, \omterm{prop:g3:sorting-living-things:groups}{oiseaux} (aussi \omterm{prop:g4:classifying-animals:classes}{reptiles}, \omterm{prop:g3:sorting-living-things:groups}{amphibiens} adultes).
\omterm{def:g7:how-animals-breathe:four}{Branchies}~: \omterm{prop:g3:sorting-living-things:groups}{poissons}, \omterm{ex:g2:animals-grow-up:frog}{têtards} (aussi moules, larves aquatiques).
\omterm{def:g7:how-animals-breathe:four}{Trachées}~: sauterelles, abeilles --- les \omterm{prop:g3:sorting-living-things:groups}{insectes} en général.
\omterm{def:g7:how-animals-breathe:four}{Respiration cutanée}~: le ver de terre~; les \omterm{prop:g3:sorting-living-things:groups}{amphibiens} en
complément.
\end{solution}

\begin{solution}{exo:g7:how-animals-breathe:3}
L'eau entre par la bouche, est poussée sur les \omterm{def:g7:how-animals-breathe:four}{branchies} plumeuses
--- où se fait l'échange --- et ressort sous les opercules~: un
courant à sens unique, pompé en alternance par la bouche et les
opercules.
\end{solution}

\begin{solution}{exo:g7:how-animals-breathe:4}
Par les \omterm{def:g7:how-animals-breathe:four}{trachées}~: un réseau de fins tubes à air partant
d'ouvertures réparties sur le corps et se ramifiant jusqu'à chaque
organe --- l'air est livré à domicile, si bien que le \omterm{prop:g4:journey-of-food:absorb}{sang} n'a
presque rien à transporter.
\end{solution}

\begin{solution}{exo:g7:how-animals-breathe:5}
Hors de l'eau, ses surfaces branchiales plumeuses s'affaissent et se
collent~: la grande surface se réduit à presque rien, et sans
surface il n'y a pas de passage --- du \omterm{def:g7:what-is-respiration:respiration}{dioxygène} en abondance,
aucune machine pour le prendre.
\end{solution}

\begin{solution}{exo:g7:how-animals-breathe:6}
La \omterm{def:g7:how-animals-breathe:four}{respiration cutanée} --- et c'est la condition d'\omterm{def:g6:exploring-our-environment:conditions}{humidité} qui
défaille~: la \omterm{def:g1:the-five-senses:sense}{peau} qui sèche ne laisse plus passer les gaz, et le
ver étouffe à l'air libre. Grande et mince ne sauvent pas une
surface qui n'est plus humide.
\end{solution}

\begin{solution}{exo:g7:how-animals-breathe:7}
Têtard~: des \omterm{def:g7:how-animals-breathe:four}{branchies}. Adulte~: des \omterm{def:g7:how-animals-breathe:four}{poumons} rudimentaires plus la
\omterm{def:g1:the-five-senses:sense}{peau} humide --- la \omterm{def:g1:the-five-senses:sense}{peau} seule assurant le service pendant l'hiver
engourdi au fond de la mare.
\end{solution}

\begin{solution}{exo:g7:how-animals-breathe:8}
Sa machine est le \omterm{def:g7:how-animals-breathe:four}{poumon} --- hérité d'ancêtres terrestres qui
allaitaient --- si bien qu'il doit remonter respirer de l'air, si
poissonnière que soit sa vie~: maintenu sous l'eau, il se noie.
L'habit est celui de l'eau~; la machine est celle de la famille.
\end{solution}

\begin{solution}{exo:g7:how-animals-breathe:9}
Une réserve d'air emportée --- une bouteille de plongée sous les
élytres, dans laquelle puisent ses \omterm{def:g7:how-animals-breathe:four}{trachées}. Il doit remonter
régulièrement en surface pour renouveler la bulle.
\end{solution}

\begin{solution}{exo:g7:how-animals-breathe:10}
(a) Moule~: résidente de l'eau, aucun mouvement visible mais des
\omterm{def:g7:how-animals-breathe:four}{branchies} à l'intérieur de la \omterm{def:g1:animals-around-us:covering}{coquille}, lavées par un courant pompé.
(b) Sauterelle~: résidente de l'air~; des rangées de minuscules
ouvertures, un abdomen qui bat --- des \omterm{def:g7:how-animals-breathe:four}{trachées}. (c) Couleuvre~: une
reptile navetteuse --- des \omterm{def:g7:how-animals-breathe:four}{poumons}~; elle nage narines hors de
l'eau et remonte respirer.
\end{solution}

\begin{solution}{exo:g7:how-animals-breathe:11}
Les \omterm{def:g7:how-animals-breathe:four}{trachées} livrent par des tubes à air ramifiés, excellentes sur
quelques millimètres et trop lentes sur de longues distances~: chez
un corps de la taille d'une baleine, les organes internes seraient
trop loin de toute ouverture. La machine fixe un plafond de taille,
et les \omterm{prop:g3:sorting-living-things:groups}{insectes} vivent en dessous.
\end{solution}

\begin{solution}{exo:g7:how-animals-breathe:12}
Le cahier des charges~: une surface de rencontre grande, mince et
humide entre le corps et le milieu. Les \omterm{def:g7:how-animals-breathe:four}{branchies} le satisfont par
des surfaces plumeuses déployées par l'eau~; les \omterm{def:g7:how-animals-breathe:four}{poumons} par des
millions de poches repliées à l'abri et humides~; les \omterm{def:g7:how-animals-breathe:four}{trachées} par
la surface intérieure totale de leurs extrémités les plus fines,
humides au bout~; la \omterm{def:g1:the-five-senses:sense}{peau} par tout l'extérieur humide du corps,
gardé grand par rapport à un corps petit et mince. Quatre supports,
une seule chose supportée.
\end{solution}

\begin{solution}{pb:g7:how-animals-breathe:1}
\textbf{1.} Alevins de truite et épinoches~: \omterm{def:g7:how-animals-breathe:four}{branchies}. \omterm{ex:g2:animals-grow-up:frog}{Têtards}~:
\omterm{def:g7:how-animals-breathe:four}{branchies} (ceux à pattes en pleine reconstruction). Grenouilles
adultes~: \omterm{def:g7:how-animals-breathe:four}{poumons} plus \omterm{def:g1:the-five-senses:sense}{peau}. Limnées~: \omterm{def:g7:how-animals-breathe:four}{poumons} --- des navetteuses
de surface. Dytique~: \omterm{def:g7:how-animals-breathe:four}{trachées}, alimentées par sa bulle emportée.
Larves d'éphémère~: \omterm{def:g7:how-animals-breathe:four}{branchies} (les touffes plumeuses). Vers de
terre~: la \omterm{def:g1:the-five-senses:sense}{peau}.

\textbf{2.} Des \omterm{def:g7:how-animals-breathe:four}{branchies}. L'agitation renouvelle l'eau au contact
de la surface branchiale --- un courant fabriqué sur place là où la
mare immobile n'en fournit aucun, qui maintient au lieu du passage
une eau fraîche chargée de \omterm{def:g7:what-is-respiration:respiration}{dioxygène}.

\textbf{3.} Les limnées, le dytique, les grenouilles adultes (et
toute couleuvre de passage)~: tous les utilisateurs de \omterm{def:g7:how-animals-breathe:four}{poumons} et de
bulles --- leurs machines prennent le \omterm{def:g7:what-is-respiration:respiration}{dioxygène} dans l'air, si bien
que la surface est leur station-service.

\textbf{4.} La reconstruction respiratoire de la \omterm{def:g2:animals-grow-up:metamorphosis}{métamorphose}~: les
\omterm{def:g7:how-animals-breathe:four}{branchies} qui servaient le têtard sans pattes sont échangées contre
les \omterm{def:g7:how-animals-breathe:four}{poumons} (et la \omterm{def:g1:the-five-senses:sense}{peau} respirante) dont le grenouillet aura besoin
à terre.

\textbf{5.} Les utilisateurs de \omterm{def:g7:how-animals-breathe:four}{branchies} souffrent en premier~:
toute leur réserve est le stock dissous qui s'épuise. Les
navetteurs --- limnée, dytique, grenouille --- puisent dans l'air
illimité du dessus et se contentent de faire la navette plus
souvent.

\textbf{6.} Le millimètre supérieur touche l'air~: le \omterm{def:g7:what-is-respiration:respiration}{dioxygène} s'y
dissout depuis le dessus, si bien que la couche de la pellicule est
l'eau la moins pauvre d'une mare qui étouffe. Les épinoches font la
queue au dernier puits.

\textbf{7.} Les \omterm{def:g7:how-animals-breathe:four}{branchies} s'encrassent~: la vase fine se dépose sur
les surfaces plumeuses et bloque le passage. Attends-toi à voir les
respirants à \omterm{def:g7:how-animals-breathe:four}{branchies} fuir le nuage, et le pompage des opercules
peiner --- des spasmes de toux pour chasser le limon.

\textbf{8.} La \omterm{def:g7:how-animals-breathe:four}{respiration cutanée}, et la condition «~humide mais
non privé d'air~»~: les galeries gorgées d'eau contiennent trop peu
de \omterm{def:g7:what-is-respiration:respiration}{dioxygène}, si bien que les vers montent respirer --- ils
échangent l'étouffement d'en bas contre le risque de dessèchement
d'en haut.

\textbf{9.} Tableau à dresser (machine et milieu pour chaque
\omterm{def:g1:animals-around-us:animal}{animal}~; navetteurs~: limnée, dytique, grenouille). En tête de
chaque colonne~: la surface d'échange grande, mince et humide --- le
cahier des charges unique que tiennent les quatre machines.

\textbf{10.} Grenouilles~: pourvues --- engourdies dans la vase,
elles respirent par la \omterm{def:g1:the-five-senses:sense}{peau} leur petit besoin d'hiver. Limnées et
dytique~: en difficulté --- leur surface est scellée~; ils hivernent
au ralenti, besoins réduits ou bulles emprisonnées. \omterm{prop:g3:sorting-living-things:groups}{Poissons}~:
pourvus tant que dure le stock de l'eau --- le froid ralentit tous
les bruits de moteur
(\cref{ex:g7:what-is-respiration:demand}), et les \omterm{def:g7:how-animals-breathe:four}{branchies}
suffisent à la demande réduite sous la glace.

\textbf{11.} La cloche de soie de l'araignée est un réservoir d'air
fixe là où la bulle du dytique est un réservoir emporté~: toutes
deux stockent de l'air de surface sous l'eau et continuent de servir
une machine à air (l'équipement de \omterm{def:g7:how-animals-breathe:four}{poumons} en feuillets et de
\omterm{def:g7:how-animals-breathe:four}{trachées} de l'araignée) dans le mauvais milieu.

\textbf{12.} Parce que la mare offre du \omterm{def:g7:what-is-respiration:respiration}{dioxygène} dans deux milieux
--- rare et dissous en dessous, abondant au-dessus --- et que chaque
machine récolte un seul milieu à une seule taille de vie~: ce n'est
qu'avec les quatre au travail que chaque recoin, de la vase à la
pellicule de surface, est habité.
\end{solution}
